\chapter{Inequalities That Count}

\begin{center}\itshape
You don't have to believe in God, but, as a mathematician, you should believe in The Book.\\
\upshape --- Paul Erdős, as quoted by M.~Aigner and G.~M.~Ziegler in \emph{Proofs from THE BOOK}
\end{center}

In 1946 Paul Erdős published a two-page note in the \emph{American Mathematical Monthly} with a question that anyone can
understand. Put $n$ points in the plane. Among the $\binom n2$ distances between them, how few different values can there be?
If the points are equally spaced on a line, there are $n-1$ distances. Erdős noticed that a square grid does much better: on a
$\sqrt n\times\sqrt n$ grid every distance is the square root of a sum of two squares below $2n$, and by a theorem of Landau
there are only about $\frac{cn}{\sqrt{\log n}}$ such numbers. He proved that one can never go below roughly $\sqrt n$ distances,
with an argument of a few lines (Problem~\ref{pb:distances} below), and he conjectured that the grid is essentially the best
possible: at least $n^{1-\varepsilon}$ distances for every $\varepsilon>0$ and $n$ large.

The gap between $\sqrt n$ and $n$ became one of the most famous open problems of combinatorial geometry. Each improvement of
the exponent $\frac12$ brought a new tool. In 1997 László Székely got $n^{4/5}$ with a ``crossing lemma'' about drawings of
graphs and a counting argument that fits on a page (its main step is Problem~\ref{pb:szemereditrotter}). Then György Elekes
had a strange idea: turn the question about distances in the plane into a question about lines in three-dimensional space,
by looking at the rigid motions that carry one pair of points onto another. Elekes died in 2008, before his programme could be
finished, and the reformulation appeared with Micha Sharir after his death. In 2010 Larry Guth and Nets Katz completed it,
with a new method that cuts space by the zero set of a well-chosen polynomial: $n$ points always determine at least
$\frac{cn}{\log n}$ distinct distances. Erdős's conjecture was settled, up to the factor $\sqrt{\log n}$ between this bound and
the grid.

At the heart of all these proofs are inequalities that count: Cauchy--Schwarz applied to the number of incidences between
points and curves, convexity applied to degrees, the entropy of a random object compared with the logarithm of the number of
choices. This chapter collects such arguments. Each problem has two numbers to compare, and the art is to find the quantity
that is counted in two ways.

\begin{think}
Seven points in the plane: how few distinct distances can they determine? Try the vertices of a regular heptagon, a regular
hexagon with its centre, and pieces of the triangular lattice. Then guess what the best configurations of many points look
like.
\end{think}

\section{Points, lines and distances}

The basic fact of plane geometry used here is that two points determine one line and two distinct circles meet in at most two
points. Turned into counting, it says that a certain bipartite graph contains no $K_{2,2}$, or no $K_{2,3}$, and the
Cauchy--Schwarz inequality does the rest (Tome~0, Theorem~4.4; compare the Kővári--Sós--Turán theorem, Tome~3, Theorem~B.6).

\begin{exercise}[$\star\star$]\label{pb:distances}
\emph{Few distances, few unit distances.} Let $P$ be a set of $n\ge2$ points in the plane.
\begin{enumerate}[label=(\alph*)]
\item (Erdős, 1946.) Show that $P$ determines at least $\sqrt{(n-2)/2}$ distinct distances.
\item Show that at most $\frac n4\bigl(1+\sqrt{8n-7}\bigr)$ pairs of points of $P$ are at distance exactly $1$.
\end{enumerate}
\end{exercise}

\begin{exercise}[$\star\star$]\label{pb:incidences}
\emph{Incidences by Cauchy--Schwarz.} Let $P$ be a set of $n$ points and $L$ a set of $m$ lines in the plane, and let $I(P,L)$ be the
number of pairs $(p,\ell)$ with $p\in\ell$.
\begin{enumerate}[label=(\alph*)]
\item Show that $I(P,L)\le n\sqrt m+m$ and $I(P,L)\le m\sqrt n+n$.
\item Let $k\ge2\sqrt n$. Show that at most $\frac{2n}k$ lines contain $k$ or more points of $P$.
\end{enumerate}
\end{exercise}

\begin{exercise}[$\star\star\star$]\label{pb:szemereditrotter}
\emph{The Szemerédi--Trotter theorem (Székely's proof).} Recall the crossing lemma (Tome~3, Exercise~12.4): a simple graph with
$v$ vertices and $e\ge4v$ edges cannot be drawn in the plane with fewer than $\frac{e^3}{64v^2}$ crossings.
\begin{enumerate}[label=(\alph*)]
\item With the notation of Problem~\ref{pb:incidences}, show that $I(P,L)\le4(nm)^{2/3}+m+4n$. (Join consecutive points of $P$ on
each line.)
\item Deduce that, for $k\ge2$, the number of lines containing at least $k$ points of $P$ is $O\bigl(\frac{n^2}{k^3}+\frac nk\bigr)$.
\item By the same method with unit circles, show that $n$ points determine at most $4n^{4/3}+5n$ unit distances.
\end{enumerate}
\end{exercise}

\begin{figure}[ht]
\centering
\fig{ch08-szekely}
\caption{Székely's graph for the points of a $4\times4$ grid and six lines: on each line, consecutive points are joined by an edge.
Two edges can cross only where their lines cross, so the graph has few crossings, while the crossing lemma says that a graph with many
edges must have many crossings (Problem~\ref{pb:szemereditrotter}).}
\label{fig:szekely}
\end{figure}

\begin{exercise}[$\star\star$]\label{pb:sumproduct}
\emph{Sums or products (Elekes, 1997).} Let $A$ be a finite set of $N$ positive real numbers. By looking at the point set
$(A+A)\times(A\cdot A)$ and the $N^2$ lines $y=b(x-c)$, $b,c\in A$, show that
\[
|A+A|\cdot|A\cdot A|\ge c\,N^{5/2}
\]
for an absolute constant $c>0$, so that $\max(|A+A|,|A\cdot A|)\ge\sqrt c\,N^{5/4}$.
\end{exercise}

\begin{figure}[ht]
\centering
\fig{ch08-elekes}
\caption{Elekes's picture for $A=\{1,\dots,5\}$: the grey grid is $(A+A)\times(A\cdot A)$, and each line $y=b(x-c)$ with $b,c\in A$ meets
it in the $5$ points $(a+c,ab)$, $a\in A$. Four of the $25$ lines are drawn. Many incidences on a small grid are forbidden by the
Szemerédi--Trotter theorem, so the grid cannot be small in both directions (Problem~\ref{pb:sumproduct}).}
\label{fig:elekes}
\end{figure}

\begin{exercise}[$\star\star$]\label{pb:energy}
\emph{Energy and the multiplication table.} For a finite set $A$ of numbers, its \emph{additive energy} $E_+(A)$ is the number of
quadruples $(a,b,c,d)\in A^4$ with $a+b=c+d$, and its multiplicative energy $E_\times(A)$ the number with $ab=cd$.
\begin{enumerate}[label=(\alph*)]
\item Show that $|A+A|\ge\frac{|A|^4}{E_+(A)}$ and $|A\cdot A|\ge\frac{|A|^4}{E_\times(A)}$.
\item Let $A=\{1,\dots,N\}$. Show that $E_\times(A)=\sum\bigl\lfloor\frac N{\max(u,v)}\bigr\rfloor^2$, the sum running over the pairs of coprime
integers $1\le u,v\le N$, and deduce $E_\times(A)=\frac{12}{\pi^2}N^2\log N+O(N^2)$.
\item Deduce that the $N\times N$ multiplication table contains at least $\bigl(\frac{\pi^2}{12}+o(1)\bigr)\frac{N^2}{\log N}$ distinct numbers.
\end{enumerate}
\end{exercise}

\section{Counting in graphs}

For a graph $G$ with adjacency matrix $A$, the number of closed walks of length $k$ is $\operatorname{tr}A^k=\sum_i\lambda_i^k$, where the
$\lambda_i$ are the (real) eigenvalues of $A$, and the number of all walks of length $k$ is $\mathbf1^{\mathsf T}A^k\mathbf1$. Inequalities on
eigenvalues become inequalities on subgraph counts (Tome~11, Appendix~A).

\begin{exercise}[$\star\star$]\label{pb:triangleseigen}
\emph{Triangles from eigenvalues.} Let $G$ be a simple graph with $m$ edges, $t$ triangles and adjacency eigenvalues
$\lambda_1\ge\dots\ge\lambda_n$.
\begin{enumerate}[label=(\alph*)]
\item Show that $\sum_i\lambda_i^2=2m$ and $\sum_i\lambda_i^3=6t$.
\item Deduce that $t\le\frac{\sqrt2}3m^{3/2}$, and that the ratio $t/m^{3/2}$ tends to $\frac{\sqrt2}3$ for complete graphs.
\end{enumerate}
\end{exercise}

\begin{exercise}[$\star\star$]\label{pb:goodman}
\emph{Goodman's bound (1959).} Colour the edges of $K_n$ red and blue.
\begin{enumerate}[label=(\alph*)]
\item Count the triangles that are not monochromatic through their two-coloured ``corners'', and show that there are at least
\[
\binom n3-\frac n2\Bigl\lfloor\frac{(n-1)^2}4\Bigr\rfloor
\]
monochromatic triangles. Check that this gives $2$ for $n=6$.
\item Show that this is asymptotically $\frac14\binom n3$, and that a random colouring has $\frac14\binom n3$ monochromatic triangles on average.
\end{enumerate}
\end{exercise}

\begin{exercise}[$\star\star$]\label{pb:walks}
\emph{Many walks, many squares.} Let $G$ be a graph with $n$ vertices, $m$ edges and average degree $\bar d=\frac{2m}n$.
\begin{enumerate}[label=(\alph*)]
\item Show that for every even $k\ge2$ the number of walks of length $k$ in $G$ is at least $n\bar d^{\,k}$.
\item Show that $\operatorname{tr}A^4\ge\bar d^{\,4}$: the number of closed walks of length $4$ is at least what a random graph with the same
density has on average, $n^4(\bar d/n)^4$, up to lower-order terms.
\item For the Paley graph $P_p$ (Problem~\ref{pb:paley}), compute $\operatorname{tr}A^4$ and show that $\operatorname{tr}A^4=\bar d^{\,4}(1+O(\frac1p))$.
\end{enumerate}
\end{exercise}

\begin{exercise}[$\star$]\label{pb:treebound}
\emph{Spanning trees and AM--GM.} Let $G$ be a connected graph with $n$ vertices and $m$ edges, and $\tau(G)$ its number of spanning
trees. Show that
\[
\tau(G)\le\frac1n\Bigl(\frac{2m}{n-1}\Bigr)^{n-1},
\]
with equality if and only if $G$ is complete. Which theorem do you recover for $K_n$?
\end{exercise}

\section{Families of sets}

A family $\mathcal F$ of subsets of $[n]=\{1,\dots,n\}$ is an \emph{antichain} if no member contains another; Sperner's theorem and the
LYM inequality bound antichains (Tome~3, Theorem~A.16). The problems below add one probabilistic idea each: a random
ordering, a random walk down the cube, or a potential function in a game.

\begin{exercise}[$\star\star$]\label{pb:bollobas}
\emph{Bollobás's set-pairs inequality (1965).} Let $A_1,\dots,A_m$ and $B_1,\dots,B_m$ be finite sets with $A_i\cap B_i=\varnothing$ for every
$i$ and $A_i\cap B_j\ne\varnothing$ for $i\ne j$. Let $a_i=|A_i|$, $b_i=|B_i|$.
\begin{enumerate}[label=(\alph*)]
\item Show that $\displaystyle\sum_{i=1}^m\binom{a_i+b_i}{a_i}^{-1}\le1$.
\item A graph is \emph{$\tau$-critical} if deleting any edge lowers the minimum size $\tau$ of a vertex cover. Show that a $\tau$-critical
graph with $\tau=t+1$ has at most $\binom{t+2}2$ edges, and that $K_{t+2}$ attains this.
\end{enumerate}
\end{exercise}

\begin{exercise}[$\star\star$]\label{pb:littlewoodofford}
\emph{The Littlewood--Offord problem (Erdős, 1945).} Let $a_1,\dots,a_n$ be real numbers with $|a_i|\ge1$, and let $\varepsilon_1,\dots,\varepsilon_n$ be
independent uniform signs $\pm1$.
\begin{enumerate}[label=(\alph*)]
\item Show that for every open interval $J$ of length $2$, $P\bigl(\sum_i\varepsilon_ia_i\in J\bigr)\le\binom n{\lfloor n/2\rfloor}2^{-n}$.
\item Show that the bound is attained, and that it is asymptotic to $\sqrt{\frac2{\pi n}}$.
\end{enumerate}
\end{exercise}

\begin{exercise}[$\star\star$]\label{pb:kleitman}
\emph{Harris and Kleitman.} A family $\mathcal A$ of subsets of $[n]$ is \emph{increasing} if $S\in\mathcal A$ and $S\subseteq T$ imply $T\in\mathcal A$,
and \emph{decreasing} if its complement in $2^{[n]}$ is increasing.
\begin{enumerate}[label=(\alph*)]
\item (Harris, 1960; Kleitman, 1966.) Show that $2^n|\mathcal A\cap\mathcal B|\ge|\mathcal A||\mathcal B|$ if $\mathcal A,\mathcal B$ are both increasing or both decreasing.
\item A family is \emph{intersecting} if any two members meet. Show that the union of $k$ intersecting families of subsets of $[n]$ has
at most $2^n-2^{n-k}$ members, and that this is attained.
\end{enumerate}
\end{exercise}

\begin{exercise}[$\star\star$]\label{pb:sauer}
\emph{The Sauer--Shelah lemma.} A family $\mathcal F$ of subsets of $[n]$ \emph{shatters} $T\subseteq[n]$ if $\{F\cap T:F\in\mathcal F\}$ is the set of
all subsets of $T$. The \emph{VC-dimension} of $\mathcal F$ is the largest size of a shattered set.
\begin{enumerate}[label=(\alph*)]
\item Show that if $\mathcal F$ has VC-dimension at most $d$, then $|\mathcal F|\le\sum_{i=0}^d\binom ni$.
\item Show that the subsets of $n$ points of the real line cut out by intervals form a family of VC-dimension $2$ that attains the
bound.
\end{enumerate}
\end{exercise}

\begin{exercise}[$\star\star$]\label{pb:cubeiso}
\emph{Edges inside a set of vertices of the cube.} Let $S$ be a set of vertices of the cube $\{0,1\}^n$, and $e(S)$ the number of edges of
the cube with both ends in $S$. Show that $e(S)\le\frac12|S|\log_2|S|$, with equality when $S$ is a subcube.
\end{exercise}

\begin{exercise}[$\star\star$]\label{pb:erdosselfridge}
\emph{The Erdős--Selfridge criterion (1973).} Maker and Breaker alternately claim unclaimed points of a finite set $V$, Maker first.
Maker wins if he claims all the points of some set of a given family $\mathcal H$ of subsets of $V$; otherwise Breaker wins.
\begin{enumerate}[label=(\alph*)]
\item Show that if $\sum_{H\in\mathcal H}2^{-|H|}<\frac12$, then Breaker has a winning strategy.
\item Maker and Breaker claim edges of the complete graph $K_N$, and Maker wants a complete subgraph $K_q$. Show that Breaker wins if
$\binom Nq2^{-\binom q2}<\frac12$, for instance if $N\le2^{q/2}$ and $q\ge3$.
\end{enumerate}
\end{exercise}

\section{Entropy and convexity}

The entropy $H(X)=\sum_xP(X=x)\log_2\frac1{P(X=x)}$ of a random variable with finitely many values is at most $\log_2$ of the number of
values (Tome~3, Chapter~4). Comparing the entropy of a uniformly random object with the entropies of its parts gives sharp
counting inequalities.

\begin{exercise}[Journey]\label{pb:shearer}
\emph{Shearer's lemma and the Loomis--Whitney inequality.} All random variables take finitely many values. For random variables
$X,Y$ define $H(Y\mid X)=\sum_xP(X=x)H(Y\mid X=x)$, where $H(Y\mid X=x)$ is the entropy of the conditional law.
\begin{enumerate}[label=(\alph*)]
\item Show the chain rule $H(X,Y)=H(X)+H(Y\mid X)$. Using subadditivity, $H(X,Y)\le H(X)+H(Y)$ (Tome~3, Exercise~4.9), deduce that
conditioning on more never raises entropy: $H(Y\mid X,Z)\le H(Y\mid X)$.
\item (Shearer, 1986.) Let $X=(X_1,\dots,X_n)$, and let $F_1,\dots,F_r\subseteq[n]$ be such that every $i\in[n]$ belongs to at least $k$ of them.
Writing $X_F=(X_i)_{i\in F}$, show that $kH(X)\le\sum_jH(X_{F_j})$.
\item (Loomis and Whitney, 1949, discrete form.) Let $S\subset\Z^3$ be finite, and $S_{xy},S_{yz},S_{xz}$ its projections on the coordinate
planes. Show that $|S|^2\le|S_{xy}|\,|S_{yz}|\,|S_{xz}|$.
\item Give a second proof of Problem~\ref{pb:triangleseigen}(b): a graph with $m$ edges has at most $\frac{(2m)^{3/2}}6$ triangles.
\end{enumerate}
\end{exercise}

\begin{exercise}[$\star$]\label{pb:chebyshevsum}
\emph{Chebyshev's sum inequality.}
\begin{enumerate}[label=(\alph*)]
\item Let $a_1\le\dots\le a_n$ and $b_1\le\dots\le b_n$. Show that $\frac1n\sum_ia_ib_i\ge\bigl(\frac1n\sum_ia_i\bigr)\bigl(\frac1n\sum_ib_i\bigr)$.
\item Let $f,g$ be non-decreasing functions and $X$ a real random variable with $f(X),g(X)$ square-integrable. Show that
$\mathbb E[f(X)g(X)]\ge\mathbb E[f(X)]\,\mathbb E[g(X)]$.
\item Deduce that the moments $M_k=\lim\frac1N\sum_{n\le N}\bigl(\frac{\varphi(n)}n\bigr)^k$ of Problem~\ref{pb:phiovern} satisfy
$M_{j+k}\ge M_jM_k$, and check it on $M_1=0.6079\ldots$, $M_2=0.4282\ldots$
\end{enumerate}
\end{exercise}

\begin{exercise}[$\star\star$]\label{pb:newton}
\emph{Newton's inequalities.} Let $P(x)=\sum_{k=0}^na_kx^k$ be a real polynomial of degree $n$ all of whose roots are real.
\begin{enumerate}[label=(\alph*)]
\item Show that for $1\le k\le n-1$,
\[
a_k^2\ge a_{k-1}a_{k+1}\Bigl(1+\frac1k\Bigr)\Bigl(1+\frac1{n-k}\Bigr).
\]
\item Deduce that the numbers $c(n,k)$ of permutations of $[n]$ with $k$ cycles form a log-concave, hence unimodal, sequence in $k$.
\end{enumerate}
\end{exercise}

\begin{exercise}[$\star$]\label{pb:bonferroni}
\emph{Bonferroni's inequalities.} Let $A_1,\dots,A_m$ be events and $S_j=\sum_{i_1<\dots<i_j}P(A_{i_1}\cap\dots\cap A_{i_j})$.
\begin{enumerate}[label=(\alph*)]
\item Show that $P(A_1\cup\dots\cup A_m)\le\sum_{j=1}^t(-1)^{j-1}S_j$ for $t$ odd, and $\ge$ for $t$ even.
\item Deduce that the probability $D_n/n!$ that a random permutation of $[n]$ has no fixed point differs from $\sum_{j=0}^t\frac{(-1)^j}{j!}$ by at
most $\frac1{(t+1)!}$, for $0\le t\le n$.
\end{enumerate}
\end{exercise}

\section{Arithmetic}

\begin{exercise}[$\star\star$]\label{pb:primitive}
\emph{Primitive sets (Erdős, 1935).} A set $A$ of integers $\ge2$ is \emph{primitive} if no element divides another. For $a\ge2$ let $P(a)$ be
its largest prime factor, and $M_a$ the set of integers $am$ whose every prime factor $q$ of $m$ satisfies $q\ge P(a)$.
\begin{enumerate}[label=(\alph*)]
\item Show that the sets $M_a$, $a\in A$, are pairwise disjoint when $A$ is primitive.
\item Show that $M_a$ has density $\frac1a\prod_{p<P(a)}\bigl(1-\frac1p\bigr)$.
\item Deduce that $\sum_{a\in A}\frac1{a\log a}\le C$ for every primitive set $A$, with an absolute constant $C$.
\end{enumerate}
\end{exercise}

\begin{exercise}[$\star\star$]\label{pb:ruzsa}
\emph{Ruzsa's inequalities.} Let $A,B,C$ be finite sets in an abelian group.
\begin{enumerate}[label=(\alph*)]
\item (Triangle inequality.) Show that $|A|\,|B-C|\le|A-B|\,|A-C|$. Deduce that if $|A+A|\le K|A|$ then $|A-A|\le K^2|A|$.
\item (Covering lemma.) Show that if $|A+B|\le K|A|$, then $B$ is covered by at most $K$ translates of $A-A$.
\end{enumerate}
\end{exercise}

\begin{exercise}[$\star\star$]\label{pb:divisorbound}
\emph{Champions of divisors.} Let $d(n)$ be the number of divisors of $n$ and $\alpha>0$. (Tome~5, Exercise~2.6, treats $\alpha=\frac12$.)
\begin{enumerate}[label=(\alph*)]
\item Show that $\max_n\frac{d(n)}{n^\alpha}$ exists and equals $\prod_p\max_{e\ge0}\frac{e+1}{p^{e\alpha}}$, a product with finitely many factors
different from $1$.
\item Let $n_\alpha$ be the smallest $n$ at which this maximum is reached, and write $n_\alpha=\prod_pp^{e_p}$. Show that $e_p$ is the least
$e\ge0$ with $\frac{e+2}{e+1}\le p^\alpha$. Deduce that $e_2\ge e_3\ge e_5\ge\cdots$, and that $p$ divides $n_\alpha$ if and only if $p^\alpha<2$.
\item Show that $d(n)\le3.53\,n^{1/3}$ for all $n$, with the maximum of $d(n)/n^{1/3}$ reached at $n=2520$, and find $n_{1/4}$.
\end{enumerate}
\end{exercise}

\begin{figure}[ht]
\centering
\fig{ch08-divisors}
\caption{The values of $d(n)/n^{1/3}$ for $n\le6000$, on a logarithmic scale in $n$ (a sample of points, and all those above $2.3$). The
maximum, $48/2520^{1/3}=3.527\ldots$, is reached at $n=2520=2^3\cdot3^2\cdot5\cdot7$; for the exponent $\frac12$ the record is at $n=12$
(Problem~\ref{pb:divisorbound}).}
\label{fig:divisors}
\end{figure}

\section{Hints}

\begin{hint}{pb:distances}
(a) Fix two points; every other point is on a circle around each. (b) The unit-distance graph has no $K_{2,3}$: count pairs of
neighbours.
\end{hint}

\begin{hint}{pb:incidences}
(a) Count pairs of points on a common line, then use $\bigl(\sum_\ell k_\ell\bigr)^2\le m\sum_\ell k_\ell^2$.
\end{hint}

\begin{hint}{pb:szemereditrotter}
(a) The graph has $I-m$ edges (at least, once lines with no points are removed) and at most $m^2$ crossings. (c) Two unit circles cross
at most twice; ignore circles with fewer than $3$ points.
\end{hint}

\begin{hint}{pb:sumproduct}
Each line contains the $N$ points $(a+c,ab)$. Apply Problem~\ref{pb:szemereditrotter}(a).
\end{hint}

\begin{hint}{pb:energy}
(a) $E_+(A)=\sum_sr(s)^2$ with $\sum_sr(s)=|A|^2$. (b) Write $a/c$ in lowest terms.
\end{hint}

\begin{hint}{pb:triangleseigen}
(b) $\sum\lambda_i^3\le\lambda_1\sum\lambda_i^2$ and $\lambda_1^2\le\sum\lambda_i^2$.
\end{hint}

\begin{hint}{pb:goodman}
A triangle that is not monochromatic has exactly two two-coloured corners. A vertex of red degree $r$ is the corner of
$r(n-1-r)$ two-coloured angles.
\end{hint}

\begin{hint}{pb:walks}
(a) Expand $\mathbf1$ in an orthonormal basis of eigenvectors and use Jensen. (b) $\lambda_1\ge\bar d$ (Rayleigh). (c) Use the counts of
common neighbours in Problem~\ref{pb:paley}(a).
\end{hint}

\begin{hint}{pb:treebound}
Matrix-tree theorem: $\tau(G)=\frac1n\mu_2\cdots\mu_n$ for the Laplacian eigenvalues.
\end{hint}

\begin{hint}{pb:bollobas}
(a) In a random ordering of the ground set, what is the chance that all of $A_i$ comes before all of $B_i$? Can it happen for two
indices? (b) Apply (a) to the edges and to the covers that miss them.
\end{hint}

\begin{hint}{pb:littlewoodofford}
Make all $a_i\ge1$, then show that the sets $\{i:\varepsilon_i=1\}$ giving sums in $J$ form an antichain.
\end{hint}

\begin{hint}{pb:kleitman}
(a) Induction on $n$, splitting according to whether $n\in S$, and Chebyshev's inequality for two terms. (b) Replace each family by
the increasing family it generates.
\end{hint}

\begin{hint}{pb:sauer}
(a) Induction on $n$: compare $\mathcal F$ with $\{F\setminus\{n\}\}$ and with the sets $F$ for which both $F$ and $F\cup\{n\}$ are in $\mathcal F$.
\end{hint}

\begin{hint}{pb:cubeiso}
Induction on $n$; split $S$ by the last coordinate. You will need $t\log_2t+2t\le(1+t)\log_2(1+t)$ on $[0,1]$.
\end{hint}

\begin{hint}{pb:erdosselfridge}
Give each set not yet touched by Breaker the weight $2^{-r}$, where $r$ is the number of its points Maker has not yet claimed, and let Breaker always reduce the
total weight as much as possible.
\end{hint}

\begin{hint}{pb:shearer}
(a) Write $\log\frac1{P(X=x,Y=y)}$ as a sum of two logarithms. (b) Chain rule, then drop part of the conditioning. (d) A random ordered triangle.
\end{hint}

\begin{hint}{pb:chebyshevsum}
$\sum_{i,j}(a_i-a_j)(b_i-b_j)\ge0$; for (b), take an independent copy of $X$.
\end{hint}

\begin{hint}{pb:newton}
Derivatives and reversals of polynomials with real roots have real roots. Reduce to a quadratic.
\end{hint}

\begin{hint}{pb:bonferroni}
A point lying in exactly $r$ of the events contributes $\sum_{j\le t}(-1)^{j-1}\binom rj=1-(-1)^t\binom{r-1}t$.
\end{hint}

\begin{hint}{pb:primitive}
(a) $a$ is made of the smallest primes of $am$. (c) Mertens: $\prod_{p<y}(1-\frac1p)\ge\frac c{\log y}$.
\end{hint}

\begin{hint}{pb:ruzsa}
(a) For each $x\in B-C$ fix one way of writing it. (b) Take a maximal set of disjoint translates $b+A$, $b\in B$.
\end{hint}

\begin{hint}{pb:divisorbound}
$d$ is multiplicative and $d(p^e)=e+1$. Compare $\frac{e+1}{p^{e\alpha}}$ with its value at $e+1$.
\end{hint}

\section{Solutions}

\begin{solution}{pb:distances}
\idea Two circles meet in at most two points, so knowing two distances to two fixed points almost determines a point. Turned
around, two points have at most two common neighbours in the unit-distance graph, and Cauchy--Schwarz bounds graphs without
$K_{2,3}$.

\writeup (a) Let $t$ be the number of distinct distances, and fix $p\ne q$ in $P$. Every other point $x$ lies on one of the $t$ circles
centred at $p$ with a radius in the distance set, and on one of the $t$ such circles centred at $q$. Two circles with distinct
centres meet in at most $2$ points, so $n-2\le2t^2$.

(b) Let $G$ be the graph on $P$ joining points at distance $1$, with degrees $d_x$ and $u$ edges. Two distinct points $y,z$ have at
most $2$ common neighbours (the intersection of two unit circles). Counting the triples $(x,\{y,z\})$ with $y,z$ neighbours of $x$,
\[
\sum_x\binom{d_x}2\le2\binom n2.
\]
By convexity (Jensen, Tome~0, Theorem~4.2), $\sum_x\binom{d_x}2\ge n\binom{2u/n}2$, so $\frac{2u}n\bigl(\frac{2u}n-1\bigr)\le2(n-1)$. Solving the
quadratic in $\frac{2u}n$ gives $\frac{2u}n\le\frac12\bigl(1+\sqrt{8n-7}\bigr)$, which is the claim.

\tools Tome~0, Theorems~4.2 and~4.4 (Jensen, Cauchy--Schwarz); Tome~3, Theorem~B.6 (the same argument for $K_{s,t}$).
\leads Erdős's own bound was $\sqrt{n-\frac34}-\frac12$. The exponent $\frac12$ in (a) was raised many times, to $\frac45$ by Székely and
$0.864$ by Katz and Tardos, before Guth and Katz's $\frac{cn}{\log n}$ (2010). For unit distances, Problem~\ref{pb:szemereditrotter}(c)
improves (b) to $O(n^{4/3})$, the best bound known; the grid gives $n^{1+c/\log\log n}$, and Erdős conjectured this is the truth.
\end{solution}

\begin{solution}{pb:incidences}
\idea Lines through a point fan out: two points lie on at most one common line. So the number of point-pairs on lines is at most
$\binom n2$, which bounds $\sum k_\ell^2$, and Cauchy--Schwarz bounds $\sum k_\ell$.

\writeup (a) Let $k_\ell=|P\cap\ell|$, so $I=\sum_\ell k_\ell$. Each pair of points lies on at most one line of $L$, so
$\sum_\ell k_\ell(k_\ell-1)\le n(n-1)$, that is $\sum_\ell k_\ell^2\le n(n-1)+I$. By Cauchy--Schwarz, $I^2\le m\sum_\ell k_\ell^2\le m(n^2+I)$. So
$I^2-mI-mn^2\le0$, and $I\le\frac{m+\sqrt{m^2+4mn^2}}2\le\frac{m+m+2n\sqrt m}2=m+n\sqrt m$. Two lines meet in at most one point, so the same
argument with the roles of points and lines exchanged gives $I\le n+m\sqrt n$.

(b) Let $L$ be the set of the $\ell$ lines with at least $k$ points. Then $k\ell\le I(P,L)\le\ell\sqrt n+n$, so $\ell(k-\sqrt n)\le n$, and since
$k-\sqrt n\ge\frac k2$, $\ell\le\frac{2n}k$.

\tools Tome~0, Theorem~4.4 (Cauchy--Schwarz).
\leads The bound $n\sqrt m+m$ is sharp for finite projective planes (Tome~8, Appendix~D), where $n=m=q^2+q+1$ and every
line has $q+1$ points. In the real plane it is not: the Szemerédi--Trotter theorem (next problem) replaces $n\sqrt m$ by $(nm)^{2/3}$,
and the difference is a topological fact, the crossing lemma, that fails over finite fields.
\end{solution}

\begin{solution}{pb:szemereditrotter}
\idea Draw a graph whose vertices are the points and whose edges are the segments between consecutive points on each line. It
has many edges if there are many incidences, but few crossings, because two lines cross once. The crossing lemma says that
these two facts are incompatible unless the incidences are few.

\writeup (a) Discard the lines containing no point of $P$; this changes neither $I$ nor the bound. On each line $\ell$ with $k_\ell\ge1$
points, join consecutive points: $k_\ell-1$ edges. Two points lie on one line at most, so the graph $G$ on $P$ is simple, with
$e=\sum_\ell(k_\ell-1)\ge I-m$ edges. It is drawn with straight segments, and two edges cross only where their lines cross, so it has
at most $\binom m2<m^2$ crossings. If $e\ge4n$, the crossing lemma gives $\frac{e^3}{64n^2}\le m^2$, so $e\le4(nm)^{2/3}$ and
$I\le4(nm)^{2/3}+m$. If $e<4n$, then $I<4n+m$. In both cases $I\le4(nm)^{2/3}+m+4n$.

(b) Let $L$ be the set of the $\ell$ lines with at least $k$ points. By (a), $k\ell\le4(n\ell)^{2/3}+\ell+4n$. If $k\ge2$, then
$k\ell-\ell\ge\frac{k\ell}2$, so $\frac{k\ell}2\le4(n\ell)^{2/3}+4n$, and one of the two terms is at least $\frac{k\ell}4$. If $\frac{k\ell}4\le4(n\ell)^{2/3}$
then $\ell\le\frac{16^3n^2}{k^3}$; otherwise $\ell\le\frac{16n}k$.

(c) Let $C$ be the set of the $n$ unit circles centred at the points of $P$. The number $u$ of unit distances satisfies
$2u=I(P,C)$, the number of pairs (point, circle through it). Circles with at most $2$ points of $P$ contribute at most $2n$ incidences.
On each circle with $k\ge3$ points, join consecutive points along the circle by arcs: $k$ arcs. Two points are joined by at most
two arcs (they lie on at most two unit circles), so after removing duplicate arcs we get a simple graph with
$e\ge\frac{I-2n}2$ edges, drawn so that two arcs cross only where their circles cross: at most $2\binom n2<n^2$ crossings. If $e\ge4n$, the
crossing lemma gives $e^3\le64n^4$, $e\le4n^{4/3}$; otherwise $e<4n$. So $I\le2e+2n\le8n^{4/3}+10n$, and $u=\frac I2\le4n^{4/3}+5n$.

\tools Tome~3, Exercise~12.4 (crossing lemma); Problem~\ref{pb:incidences}.
\leads Szemerédi and Trotter proved (a) in 1983 with a much longer argument; Székely's proof appeared in 1997. The theorem is sharp
up to the constant: take the points of an $N\times2N^2$ grid and the lines $y=ax+b$ with $a\le N$, $b\le N^2$. Over $\mathbb F_q$ the
crossing lemma fails, and incidence bounds there are a current research topic.
\end{solution}

\begin{solution}{pb:sumproduct}
\idea If $A$ has few sums and few products, the grid $(A+A)\times(A\cdot A)$ is small, yet it carries $N^2$ lines with $N$ points each.
That is too many incidences for Szemerédi--Trotter.

\writeup Let $P=(A+A)\times(A\cdot A)$, with $|P|=|A+A|\,|A\cdot A|$, and for $b,c\in A$ let $\ell_{b,c}$ be the line $y=b(x-c)$. These $N^2$
lines are distinct ($b$ is the slope, $c$ the root). For $a\in A$, the point $(a+c,ab)$ lies on $\ell_{b,c}$ and in $P$, and these $N$ points
are distinct. So $I(P,L)\ge N^3$, and Problem~\ref{pb:szemereditrotter}(a) gives
\[
N^3\le4\bigl(|P|N^2\bigr)^{2/3}+N^2+4|P|.
\]
If $4|P|\ge\frac{N^3}3$, then $|P|\ge\frac{N^3}{12}\ge\frac{N^{5/2}}{12}$. Otherwise, for $N\ge3$, $N^2\le\frac{N^3}3$, so
$4(|P|N^2)^{2/3}\ge\frac{N^3}3$, that is $|P|\ge\frac{N^{5/2}}{24^{3/2}}$. For $N\le2$ the inequality $|P|\ge1\ge cN^{5/2}$ holds with $c=2^{-5/2}$.
So $|P|\ge cN^{5/2}$ with $c=24^{-3/2}$, and the larger of the two factors is at least $\sqrt c\,N^{5/4}$.

\tools Problem~\ref{pb:szemereditrotter}.
\leads Erdős and Szemerédi conjectured (1983) that $\max(|A+A|,|A\cdot A|)\ge N^{2-\varepsilon}$ for integers. Elekes's $\frac54$ was improved by
Solymosi to $\frac43$ (up to logarithms), by a clever count of slopes, and later slightly beyond $\frac43$ by Konyagin and Shkredov
and others. An arithmetic progression has few sums and a geometric progression few products: no set has both.
\end{solution}

\begin{solution}{pb:energy}
\idea The energy counts collisions of sums. If there are few distinct sums, collisions are forced; Cauchy--Schwarz turns this into a
bound. For products of integers up to $N$, the solutions of $ab=cd$ can be listed explicitly.

\writeup (a) Let $r(s)$ be the number of pairs $(a,b)\in A^2$ with $a+b=s$. Then $\sum_sr(s)=|A|^2$ and $\sum_sr(s)^2=E_+(A)$. By
Cauchy--Schwarz over the $|A+A|$ values of $s$, $|A|^4\le|A+A|\,E_+(A)$. The multiplicative case is identical.

(b) Let $ab=cd$ with $a,b,c,d\in[1,N]$. Let $g=\gcd(a,c)$, $a=gu$, $c=gv$ with $u,v$ coprime. Then $ub=vd$, so $v\mid b$; write $b=lv$, and then
$d=lu$. Conversely every choice of coprime $u,v$ and integers $g,l\ge1$ gives a solution $(gu,lv,gv,lu)$, which lies in $[1,N]^4$ if and only
if $g,l\le\frac N{\max(u,v)}$; and $(u,v,g,l)$ is recovered from the solution as $g=\gcd(a,c)$. This proves the formula. For $M\ge2$ there
are $2\varphi(M)$ coprime pairs with $\max(u,v)=M$, and one pair, $(1,1)$, for $M=1$. Since $\lfloor N/M\rfloor^2=\frac{N^2}{M^2}+O(\frac NM)$,
\[
\begin{aligned}
E_\times&=N^2+2N^2\sum_{2\le M\le N}\frac{\varphi(M)}{M^2}\\
&\qquad+O\Bigl(N\sum_{M\le N}\frac{\varphi(M)}M\Bigr)\\
&=2N^2\sum_{M\le N}\frac{\varphi(M)}{M^2}+O(N^2).
\end{aligned}
\]
By partial summation from $\Phi(x)=\sum_{M\le x}\varphi(M)$, which is $\frac3{\pi^2}x^2+O(x\log x)$ (Tome~5, Theorem~2.6),
\[
\begin{aligned}
\sum_{M\le N}\frac{\varphi(M)}{M^2}&=\frac{\Phi(N)}{N^2}+2\int_1^N\frac{\Phi(x)}{x^3}dx\\
&=\frac6{\pi^2}\log N+O(1),
\end{aligned}
\]
so $E_\times=\frac{12}{\pi^2}N^2\log N+O(N^2)$.

(c) By (a) and (b), $|A\cdot A|\ge\frac{N^4}{E_\times}=\frac{\pi^2}{12}\frac{N^2}{\log N}(1+o(1))$.

\tools Tome~0, Theorem~4.4 (Cauchy--Schwarz); Tome~5, Theorem~2.6 (the mean of $\varphi$); Tome~1, Chapter~1 (gcd).
\leads Tome~5 shows that the table has $o(N^2)$ entries; the truth, found by Ford in 2008, is of order
$\frac{N^2}{(\log N)^{\delta}(\log\log N)^{3/2}}$ with $\delta=1-\frac{1+\log\log2}{\log2}=0.086\ldots$, so the energy bound is off only by a power of
$\log N$. Direct computation gives both.
\end{solution}

\begin{solution}{pb:triangleseigen}
\idea Closed walks of length $3$ are triangles traversed in one of six ways, and the eigenvalues add up their counts. The third
moment of the eigenvalues is bounded by the largest eigenvalue times the second moment.

\writeup (a) $A$ is real symmetric, so $\operatorname{tr}A^k=\sum_i\lambda_i^k$. Now $(A^2)_{vv}=d_v$, so $\operatorname{tr}A^2=\sum_vd_v=2m$; and $(A^3)_{vv}$ counts closed
walks $v\to x\to y\to v$, that is ordered pairs $(x,y)$ such that $vxy$ is a triangle, so every triangle is counted $6$ times in
$\operatorname{tr}A^3$.

(b) $6t=\sum\lambda_i^3\le\sum\lambda_i^2\,|\lambda_i|\le\lambda_{\max}\sum\lambda_i^2$, where $\lambda_{\max}=\max|\lambda_i|$, and
$\lambda_{\max}^2\le\sum\lambda_i^2=2m$. So $6t\le(2m)^{3/2}$, that is $t\le\frac{\sqrt2}3m^{3/2}$. For $K_q$, $t=\binom q3\sim\frac{q^3}6$ and
$m=\binom q2\sim\frac{q^2}2$, so $\frac t{m^{3/2}}\to\frac{1/6}{(1/2)^{3/2}}=\frac{\sqrt2}3$.

\tools Tome~11, Appendix~A (adjacency eigenvalues); Tome~0, Chapter~7 (spectral theorem).
\leads The exact maximum number of triangles with $m$ edges is given by the Kruskal--Katona theorem, and it is attained by
taking edges as ``close to a clique'' as possible. Problem~\ref{pb:shearer}(d) proves the same bound with entropy, and
Problem~\ref{pb:goodman} gives a lower bound for two colours.
\end{solution}

\begin{solution}{pb:goodman}
\idea A triangle that is not monochromatic has two corners where the colours change. Counting such corners vertex by vertex
gives at most $\frac{(n-1)^2}4$ per vertex.

\writeup (a) Call an \emph{angle} a pair of edges at a common vertex, and say it is two-coloured if its edges have different colours. A
monochromatic triangle has no two-coloured angle, and every other triangle has exactly two (with colours $RRB$, the two angles at
the ends of the blue edge). Every angle lies in exactly one triangle. So if $T$ is the number of monochromatic triangles,
\[
2\Bigl(\binom n3-T\Bigr)=\sum_vr_v(n-1-r_v),
\]
where $r_v$ is the red degree of $v$. Since $r(n-1-r)\le\lfloor\frac{(n-1)^2}4\rfloor$ for integers $r$, $T\ge\binom n3-\frac n2\lfloor\frac{(n-1)^2}4\rfloor$.
For $n=6$: $20-3\cdot6=2$.

(b) The bound is $\frac{n^3}6-\frac{n^3}8+O(n^2)=\frac{n^3}{24}+O(n^2)\sim\frac14\binom n3$. In a uniformly random colouring, each triangle is
monochromatic with probability $2\cdot2^{-3}=\frac14$, and linearity of expectation gives $\frac14\binom n3$.

\tools Tome~0, Chapter~8 (double counting); Tome~3, Chapter~2 (linearity of expectation).
\leads Since $R(3,3)=6$ (Tome~3, Appendix~B), every colouring of $K_6$ has a monochromatic triangle; Goodman's bound says
it has at least two, which is sharp. That the random colouring is asymptotically optimal is the first case of a
conjecture of Erdős on $K_q$, which turned out to be false already for $K_4$ (Thomason, 1989).
\end{solution}

\begin{solution}{pb:walks}
\idea Walks are moments of the spectral measure seen from the uniform vector, and convexity compares a moment with a power of the
first moment. For closed $4$-walks, the top eigenvalue alone suffices.

\writeup (a) Let $(u_i)$ be an orthonormal basis of eigenvectors, $A u_i=\lambda_iu_i$, and $\mathbf1=\sum_ic_iu_i$. Then $\sum c_i^2=n$,
$\mathbf1^{\mathsf T}A\mathbf1=\sum c_i^2\lambda_i=2m$, and the number of walks of length $k$ is $w_k=\mathbf1^{\mathsf T}A^k\mathbf1=\sum c_i^2\lambda_i^k$. For even
$k$, $x\mapsto x^k$ is convex on $\R$, and Jensen with the weights $\frac{c_i^2}n$ gives $\frac{w_k}n\ge\bigl(\frac{2m}n\bigr)^k$.

(b) By Rayleigh's principle (Tome~11, Lemma~A.1), $\lambda_1\ge\frac{\mathbf1^{\mathsf T}A\mathbf1}{\mathbf1^{\mathsf T}\mathbf1}=\bar d$. So
$\operatorname{tr}A^4=\sum\lambda_i^4\ge\lambda_1^4\ge\bar d^{\,4}$.

(c) In $P_p$, $\bar d=d=\frac{p-1}2$, and $(A^2)_{xy}$ is $d$ if $x=y$, $\frac{p-5}4$ if $x,y$ are adjacent and $\frac{p-1}4$ otherwise
(Problem~\ref{pb:paley}(a)). So
\[
\operatorname{tr}A^4=\sum_{x,y}(A^2)_{xy}^2=pd^2+pd\Bigl(\frac{p-5}4\Bigr)^2+pd\Bigl(\frac{p-1}4\Bigr)^2.
\]
Here $pd^2=O(p^3)$, and each of the last two terms is $\frac p2\cdot p\cdot\frac{p^2}{16}+O(p^3)=\frac{p^4}{32}+O(p^3)$, so together they
give $\frac{p^4}{16}+O(p^3)$. Since $d^4=\frac{p^4}{16}+O(p^3)$ as well, $\operatorname{tr}A^4=d^4+O(p^3)=d^4(1+O(\frac1p))$.

\tools Tome~0, Chapter~7 (spectral theorem), Theorem~4.2 (Jensen); Tome~11, Lemma~A.1; Problem~\ref{pb:paley}.
\leads (a) holds for odd $k$ too (Blakley and Roy, 1965), but the proof is harder. (b) is the case of the $4$-cycle in Sidorenko's
conjecture, still open in general. Chung, Graham and Wilson (1989) showed that a sequence of graphs of density $\frac12$ with
$\operatorname{tr}A^4=(\frac n2)^4(1+o(1))$ behaves like a random graph in many other ways (``quasi-randomness''); the Paley graphs are the
standard example.
\end{solution}

\begin{solution}{pb:treebound}
\idea The matrix-tree theorem writes $\tau(G)$ as a product of eigenvalues whose sum is known; AM--GM bounds a product with a given
sum.

\writeup Let $0=\mu_1<\mu_2\le\dots\le\mu_n$ be the Laplacian eigenvalues ($\mu_2>0$ as $G$ is connected). By the matrix-tree theorem
(Tome~11, Corollary~A.5), $\tau(G)=\frac1n\mu_2\cdots\mu_n$, and $\sum_i\mu_i=\operatorname{tr}L=\sum_vd_v=2m$. By AM--GM (Tome~0,
Theorem~4.3),
\[
\mu_2\cdots\mu_n\le\Bigl(\frac{\mu_2+\dots+\mu_n}{n-1}\Bigr)^{n-1}=\Bigl(\frac{2m}{n-1}\Bigr)^{n-1},
\]
with equality if and only if $\mu_2=\dots=\mu_n=\mu$. In that case $L=\mu(I-\frac1nJ)$ (both sides vanish on $\mathbf1$ and are $\mu$ times the
identity on $\mathbf1^\perp$), so every off-diagonal entry of $L$ is $-\frac\mu n\ne0$: all edges are present, $G=K_n$. For $K_n$,
$m=\binom n2$ and the bound is $\frac1nn^{n-1}=n^{n-2}$: Cayley's formula.

\tools Tome~11, Corollary~A.5 (matrix-tree theorem); Tome~0, Theorem~4.3 (AM--GM).
\leads For $d$-regular graphs this gives $\tau\le\frac1n\bigl(\frac{dn}{n-1}\bigr)^{n-1}$, roughly $d^n$; McKay (1983) found the sharp exponential rate
for large $d$-regular graphs, $\tau^{1/n}\to\frac{(d-1)^{d-1}}{(d^2-2d)^{d/2-1}}$, attained by random regular graphs.
\end{solution}

\begin{solution}{pb:bollobas}
\idea Order the ground set at random. The event ``all of $A_i$ before all of $B_i$'' has probability $1/\binom{a_i+b_i}{a_i}$, and two such
events cannot happen together, because of an element of $A_i\cap B_j$ and one of $A_j\cap B_i$.

\writeup (a) Let $X$ be the union of all the sets, and $\sigma$ a uniformly random ordering of $X$. Let $E_i$ be the event that every element
of $A_i$ comes before every element of $B_i$. Only the relative order of the $a_i+b_i$ elements of $A_i\cup B_i$ matters, and all orders
are equally likely; exactly $a_i!\,b_i!$ of them put $A_i$ first. So $P(E_i)=\frac{a_i!b_i!}{(a_i+b_i)!}$. Suppose $E_i$ and $E_j$ both occur, $i\ne j$.
Take $x\in A_i\cap B_j$ and $y\in A_j\cap B_i$. By $E_i$, $x$ (in $A_i$) comes before $y$ (in $B_i$); by $E_j$, $y$ (in $A_j$) comes before $x$ (in
$B_j$). This is absurd, so the events are disjoint, and $\sum_iP(E_i)\le1$.

(b) Let $G$ be $\tau$-critical with $\tau(G)=t+1$ and edges $e_1,\dots,e_m$. For each $i$, $G-e_i$ has a vertex cover $B_i$ of size $t$. It
does not cover $e_i$ (otherwise it would cover $G$), so $B_i\cap e_i=\varnothing$, and it covers every $e_j$, $j\ne i$. With $A_i=e_i$, (a) gives
$m\binom{t+2}2^{-1}\le1$. In $K_{t+2}$, a vertex cover misses at most one vertex, so $\tau=t+1$, and deleting the edge $uv$ allows the cover
$V\setminus\{u,v\}$ of size $t$.

\tools Tome~3, Chapter~1 (uniform random permutations); Tome~3, Theorem~A.16 (LYM, which is the case $B_i=X\setminus A_i$).
\leads The inequality has a ``skew'' version, where one only assumes $A_i\cap B_j\ne\varnothing$ for $i<j$; Lovász proved it with exterior
algebra (Tome~2), a striking use of linear algebra in combinatorics. Bollobás's inequality is a basic tool in the theory of
saturated graphs and hypergraphs.
\end{solution}

\begin{solution}{pb:littlewoodofford}
\idea Make all coefficients positive. Then raising one more sign from $-1$ to $+1$ increases the sum by at least $2$, so the sets of
indices with sum in a window of length $2$ form an antichain, and Sperner applies.

\writeup (a) Replacing $a_i$ by $-a_i$ and $\varepsilon_i$ by $-\varepsilon_i$ does not change the law of the sum, so we may assume $a_i\ge1$. For $S\subseteq[n]$ let
$\sigma(S)=\sum_{i\in S}a_i-\sum_{i\notin S}a_i$, the sum when $\varepsilon_i=1$ exactly for $i\in S$. If $S\subsetneq T$, then $\sigma(T)-\sigma(S)=2\sum_{i\in T\setminus S}a_i\ge2$, so
$\sigma(S)$ and $\sigma(T)$ are not both in the open interval $J$. Thus $\{S:\sigma(S)\in J\}$ is an antichain, of size at most $\binom n{\lfloor n/2\rfloor}$
by Sperner's theorem, and each $S$ has probability $2^{-n}$.

(b) With all $a_i=1$, $\sum\varepsilon_ia_i=2|S|-n$, and the interval
$J=]\,2\lfloor\frac n2\rfloor-n-1,\,2\lfloor\frac n2\rfloor-n+1[$ contains exactly the value $2\lfloor\frac n2\rfloor-n$, taken for the $\binom n{\lfloor n/2\rfloor}$
sets with $|S|=\lfloor\frac n2\rfloor$. By Stirling's formula, $\binom{2k}k2^{-2k}\sim\frac1{\sqrt{\pi k}}$, which is $\sqrt{\frac2{\pi n}}$ for $n=2k$, and the odd
case is the same since $\binom{2k+1}k2^{-2k-1}=\binom{2k}k2^{-2k}\cdot\frac{2k+1}{2k+2}$.

\tools Tome~3, Theorem~A.16 (LYM and Sperner); Tome~0, Exercise~5.7 (Stirling's formula).
\leads Littlewood and Offord met this question in 1943 while counting real roots of random polynomials. For coefficients that are
also distinct, Erdős and Moser conjectured, and Stanley proved with the hard Lefschetz theorem, a bound of order $n^{-3/2}$.
Tao and Vu's ``inverse'' theorems explain which $a_i$ give large concentration, and they are central in random matrix theory.
\end{solution}

\begin{solution}{pb:kleitman}
\idea Split the cube into two halves according to the last coordinate. Each half of an increasing family is increasing, and the
upper half contains the lower one; Chebyshev's inequality for two numbers glues the halves.

\writeup (a) Induction on $n$; for $n=0$ the families are $\varnothing$ or $\{\varnothing\}$ and the claim is clear. For $n\ge1$ and an increasing $\mathcal A$, let
$\mathcal A_0=\{S\subseteq[n-1]:S\in\mathcal A\}$ and $\mathcal A_1=\{S\subseteq[n-1]:S\cup\{n\}\in\mathcal A\}$. Both are increasing families on $[n-1]$, and $\mathcal A_0\subseteq\mathcal A_1$. Do
the same for $\mathcal B$. Then, by induction,
\[
|\mathcal A\cap\mathcal B|=|\mathcal A_0\cap\mathcal B_0|+|\mathcal A_1\cap\mathcal B_1|\ge\frac{|\mathcal A_0||\mathcal B_0|+|\mathcal A_1||\mathcal B_1|}{2^{n-1}}\ge\frac{(|\mathcal A_0|+|\mathcal A_1|)(|\mathcal B_0|+|\mathcal B_1|)}{2^n}=\frac{|\mathcal A||\mathcal B|}{2^n},
\]
where the last inequality is $2(xy+XY)-(x+X)(y+Y)=(X-x)(Y-y)\ge0$ with $x=|\mathcal A_0|\le X=|\mathcal A_1|$ and $y\le Y$ likewise. For decreasing
families apply this to the increasing families $\{[n]\setminus S:S\in\mathcal A\}$.

(b) If $\mathcal F$ is intersecting, so is its up-closure $\{T:T\supseteq S\text{ for some }S\in\mathcal F\}$, which contains $\mathcal F$; so we may assume the
families $\mathcal F_1,\dots,\mathcal F_k$ increasing. An intersecting family contains at most one of $S$ and $[n]\setminus S$, so $|\mathcal F_j|\le2^{n-1}$, and its complement
$\mathcal G_j$ in $2^{[n]}$ is decreasing with $|\mathcal G_j|\ge2^{n-1}$. By (a) applied repeatedly (an intersection of decreasing families is decreasing),
$|\mathcal G_1\cap\dots\cap\mathcal G_k|\ge2^n\prod_j\frac{|\mathcal G_j|}{2^n}\ge2^{n-k}$. The union of the $\mathcal F_j$ is the complement of this intersection. The families
$\mathcal F_j=\{S:j\in S\}$, $j\le k$, are intersecting, and their union misses exactly the $2^{n-k}$ sets disjoint from $[k]$.

\tools Tome~0, Chapter~8 (counting); Problem~\ref{pb:chebyshevsum} (the two-term case).
\leads (a) is the simplest case of the FKG inequality (Fortuin, Kasteleyn and Ginibre, 1971), central in statistical physics: in a
random graph, increasing properties such as ``connected'' and ``contains a triangle'' are positively correlated. Kleitman's (b)
was the answer to a question of Erdős.
\end{solution}

\begin{solution}{pb:sauer}
\idea Delete the last point. The family shrinks only where two sets differ just by that point, and those pairs form a family of
smaller VC-dimension.

\writeup (a) Let $\Phi(n,d)=\sum_{i\le d}\binom ni$, with $\Phi(n,d)=2^n$ for $d\ge n$ and $\Phi(n,-1)=0$; Pascal's rule gives $\Phi(n,d)=\Phi(n-1,d)+\Phi(n-1,d-1)$.
Induction on $n$, the case $n=0$ being clear. Let $\mathcal F'=\{F\setminus\{n\}:F\in\mathcal F\}$ and $\mathcal F''=\{F\subseteq[n-1]:F\in\mathcal F\text{ and }F\cup\{n\}\in\mathcal F\}$. The
map $F\mapsto F\setminus\{n\}$ from $\mathcal F$ to $\mathcal F'$ is two-to-one exactly above $\mathcal F''$, so $|\mathcal F|=|\mathcal F'|+|\mathcal F''|$. A set shattered by $\mathcal F'$ is shattered by $\mathcal F$,
so $\mathcal F'$ has VC-dimension $\le d$. If $\mathcal F''$ shatters $T\subseteq[n-1]$, then $\mathcal F$ shatters $T\cup\{n\}$, so $\mathcal F''$ has VC-dimension $\le d-1$. By
induction $|\mathcal F|\le\Phi(n-1,d)+\Phi(n-1,d-1)=\Phi(n,d)$.

(b) Let the points be $x_1<\dots<x_n$. An interval cuts out $\varnothing$ or a block $\{x_i,\dots,x_j\}$, $i\le j$: $1+n+\binom n2=\Phi(n,2)$ sets. Two points are
shattered, but among three points $x_i<x_j<x_k$ no interval cuts out $\{x_i,x_k\}$, so the VC-dimension is $2$.

\tools Tome~0, Chapter~1 (Pascal's rule).
\leads The lemma was found independently by Sauer, by Shelah, and by Vapnik and Chervonenkis around 1971--72. In learning
theory it is the reason why finite VC-dimension implies that empirical averages converge uniformly, and in discrete geometry
it gives the $\varepsilon$-net theorem of Haussler and Welzl.
\end{solution}

\begin{solution}{pb:cubeiso}
\idea Split the cube along one coordinate. Edges inside $S$ are edges inside each half plus a matching between the halves, at most
the size of the smaller half; an entropy-like inequality closes the induction.

\writeup Write $\log=\log_2$ and $0\log0=0$. Induction on $n$; for $n=0$, $e(S)=0$. Let $S_0,S_1$ be the parts of $S$ with last coordinate $0$
and $1$, seen in $\{0,1\}^{n-1}$, with $|S_0|=a\ge|S_1|=b$ (or the reverse). Edges inside $S$ are edges inside $S_0$, inside $S_1$, and edges in
the last direction, which join $x0$ to $x1$ and number at most $b$. So, by induction, $e(S)\le\frac12(a\log a+b\log b)+b$, and it suffices
that $a\log a+b\log b+2b\le(a+b)\log(a+b)$. For $b=0$ this is clear; otherwise put $t=\frac ba\in]0,1]$ and divide by $a$: it becomes
$g(t)=(1+t)\log(1+t)-t\log t-2t\ge0$. Now $g(0^+)=0$, $g(1)=2-2=0$, and $g''(t)=\frac1{\ln2}\bigl(\frac1{1+t}-\frac1t\bigr)<0$: $g$ is concave on $[0,1]$
and vanishes at the ends, so $g\ge0$.

A subcube of dimension $k$ has $2^k$ vertices and $k2^{k-1}=\frac12\cdot2^k\cdot k$ edges.

\tools Tome~0, Chapter~4 (convexity).
\leads This edge-isoperimetric inequality (Harper 1964, Bernstein 1967, Hart 1976) says equivalently that a set of $2^k$ vertices has at least
$2^k(n-k)$ edges leaving it. It is the combinatorial shadow of the logarithmic Sobolev inequality on the cube, and it underlies
the analysis of Boolean functions (Kahn--Kalai--Linial).
\end{solution}

\begin{solution}{pb:erdosselfridge}
\idea A potential function measures how dangerous the position is for Breaker: the sum over live sets of $2^{-(\text{points Maker still
needs})}$. A completed set has potential $1$. Maker's move can only add what Breaker's greedy move just removed.

\writeup (a) At any moment, call a set $H\in\mathcal H$ \emph{alive} if Breaker has no point of $H$, give it weight $w(H)=2^{-u(H)}$ with $u(H)$ the number of
points of $H$ not claimed by Maker, and let $\Phi=\sum_{H\text{ alive}}w(H)$. Initially $\Phi=\sum2^{-|H|}<\frac12$. Maker's first move at a point $x$ doubles the
weights of the alive sets containing $x$, so afterwards $\Phi<1$.

Consider a round in which Breaker claims $y$ and then Maker claims $x$. Breaker chooses $y$ maximising $W(y)=\sum_{H\text{ alive},\,H\ni y}w(H)$ among
the free points. His move kills those sets: $\Phi$ decreases by $W(y)$. Maker's move then doubles the weights of the alive sets containing
$x$ (none contains $y$ now), increasing $\Phi$ by at most $W(x)$, computed before Breaker's move, and $W(x)\le W(y)$ by the choice of $y$. So $\Phi$
does not increase over a round, and stays $<1$ until the end. If Maker ever claimed all of some $H$, that set would be alive with
weight $2^0=1\le\Phi$. So Maker never completes a set.

(b) Take $V$ the edge set of $K_N$ and $\mathcal H$ the edge sets of the $\binom Nq$ copies of $K_q$, each of size $\binom q2$. If $N\le2^{q/2}$, then
$\binom Nq2^{-\binom q2}\le\frac{N^q}{q!}2^{-q(q-1)/2}\le\frac{2^{q/2}}{q!}<\frac12$ for $q\ge3$.

\tools Tome~3, Theorem~12.2 (the probabilistic method, whose ``derandomisation'' this is); Tome~0, Chapter~8.
\leads The potential is the conditional expectation, under random play from now on, of the number of sets Maker would complete,
and the proof is the method of conditional expectations. In the clique game the Erd\H os--Selfridge bound is sharp in the exponent: Beck (2002) showed that the largest clique Maker can
build on $K_N$ has $2\log_2N-2\log_2\log_2N+O(1)$ vertices, so Maker wins the $K_q$ game once $N\ge2^{q/2+o(q)}$.
\end{solution}

\begin{solution}{pb:shearer}
\idea Entropy is additive along a chain and decreases when we condition on more. In Shearer's inequality each coordinate is
counted $k$ times on the right, each time with less conditioning than on the left.

\writeup (a) Let $p(x)=P(X=x)$ and $p(y\mid x)=P(Y=y\mid X=x)$. Then $\log\frac1{p(x)p(y|x)}=\log\frac1{p(x)}+\log\frac1{p(y|x)}$; averaging over $(X,Y)$ gives
$H(X,Y)=H(X)+H(Y\mid X)$. With subadditivity this gives $H(Y\mid X)=H(X,Y)-H(X)\le H(Y)$: conditioning on one variable never
raises entropy. Apply this to the conditional law of $(Z,Y)$ given $X=x$: it gives $H(Y\mid X=x,Z)\le H(Y\mid X=x)$, and averaging over
$x$ with the weights $P(X=x)$ gives $H(Y\mid X,Z)\le H(Y\mid X)$. Finally, Gibbs's inequality
against the uniform law gives $H(X)\le\log_2$(number of values).

(b) Iterating (a), $H(X)=\sum_iH(X_i\mid X_1,\dots,X_{i-1})$ (chain rule), and similarly, listing $F_j$ in increasing order,
\[
H(X_{F_j})=\sum_{i\in F_j}H\bigl(X_i\mid X_{i'},\ i'\in F_j,\ i'<i\bigr)\ge\sum_{i\in F_j}H(X_i\mid X_1,\dots,X_{i-1}),
\]
by (a), since $\{i'\in F_j:i'<i\}\subseteq\{1,\dots,i-1\}$. Summing over $j$, each $i$ appears at least $k$ times and the terms are non-negative, so
$\sum_jH(X_{F_j})\ge k\sum_iH(X_i\mid X_{<i})=kH(X)$.

(c) Let $X=(X_1,X_2,X_3)$ be uniform on $S$, so $H(X)=\log|S|$. The sets $\{1,2\},\{2,3\},\{1,3\}$ cover each coordinate twice, and $X_{\{1,2\}}$ takes
values in $S_{xy}$, so $H(X_{\{1,2\}})\le\log|S_{xy}|$, and similarly. By (b), $2\log|S|\le\log|S_{xy}|+\log|S_{yz}|+\log|S_{xz}|$.

(d) Let $(X_1,X_2,X_3)$ be uniform among the $6t$ ordered triples of vertices forming a triangle. Each pair $(X_i,X_j)$ is an ordered edge, one of
$2m$ values. By (b) with the same three pairs, $2\log(6t)\le3\log(2m)$, that is $6t\le(2m)^{3/2}$.

\tools Tome~3, Exercise~4.9 (subadditivity), Proposition~4.16 (Gibbs's inequality) and Chapter~4 (entropy).
\leads Shearer's lemma is due to Chung, Frankl, Graham and Shearer (1986). The same method proves Brègman's bound on permanents
(Radhakrishnan, 1997) and Kahn's bound on independent sets in regular bipartite graphs. Loomis and Whitney proved their inequality
for volumes of bodies in $\R^n$; it implies the isoperimetric inequality up to a constant.
\end{solution}

\begin{solution}{pb:chebyshevsum}
\idea Two similarly ordered sequences go up together, so every product $(a_i-a_j)(b_i-b_j)$ is non-negative; summing them is exactly the
inequality.

\writeup (a) $\sum_{i,j}(a_i-a_j)(b_i-b_j)\ge0$ since each term is a product of two numbers of the same sign. Expanding,
$2n\sum_ia_ib_i-2\sum_ia_i\sum_jb_j\ge0$.

(b) Let $Y$ be an independent copy of $X$. Then $(f(X)-f(Y))(g(X)-g(Y))\ge0$, and taking expectations,
$2\mathbb E[fg(X)]-2\mathbb E[f(X)]\mathbb E[g(X)]\ge0$.

(c) $M_k=\mathbb E[R^k]$ for a random variable $R$ with values in $[0,1]$: the limiting law of $\varphi(n)/n$ (Problem~\ref{pb:phiovern}). The functions
$x^j$ and $x^k$ are non-decreasing on $[0,1]$, so by (b) $M_{j+k}\ge M_jM_k$. Numerically $M_2=0.4282\ge M_1^2=0.3696$.

\tools Tome~3, Chapter~5 (expectation, independence); Problem~\ref{pb:phiovern}.
\leads (b) is the one-dimensional FKG inequality; Problem~\ref{pb:kleitman} is its version on the cube. Chebyshev used (a) in
1882 for integrals; with $f=g$ it is just $\operatorname{Var}f(X)\ge0$.
\end{solution}

\begin{solution}{pb:newton}
\idea Differentiating and reversing a real-rooted polynomial keeps it real-rooted (Rolle). Doing this cleverly reduces any three
consecutive coefficients to a quadratic, whose discriminant must be non-negative.

\writeup (a) Say a polynomial \emph{splits} if it is a product of real linear factors (constants included). If $Q$ splits, so does $Q'$:
between consecutive distinct roots there is a root of $Q'$ by Rolle, a root of multiplicity $r$ is a root of $Q'$ of multiplicity
$r-1$, and this accounts for $\deg Q-1$ roots. If $Q$ splits and has degree at most $N$, so does $x^NQ(1/x)$, whose non-zero roots are the
inverses of those of $Q$.

Let $Q=P^{(k-1)}=\sum_{j\ge k-1}a_j\frac{j!}{(j-k+1)!}x^{j-k+1}$, of degree $n-k+1$, and $R(x)=x^{n-k+1}Q(1/x)=\sum_{i=0}^{n-k+1}a_{i+k-1}\frac{(i+k-1)!}{i!}x^{n-k+1-i}$.
Differentiating $R$ $(n-k-1)$ times kills all terms with $i\ge3$ and leaves the split polynomial
\[
a_{k-1}(k-1)!\frac{(n-k+1)!}2x^2+a_kk!\,(n-k)!\,x+a_{k+1}\frac{(k+1)!}2(n-k-1)!.
\]
If its leading coefficient is $0$ the inequality reads $a_k^2\ge0$. Otherwise a split quadratic has a non-negative discriminant:
\[
a_k^2k!^2(n-k)!^2\ge a_{k-1}a_{k+1}(k-1)!(k+1)!(n-k+1)!(n-k-1)!,
\]
which is $a_k^2\ge a_{k-1}a_{k+1}\frac{k+1}k\cdot\frac{n-k+1}{n-k}$.

(b) $\sum_kc(n,k)x^k=x(x+1)\cdots(x+n-1)$ (the cycle decomposition, built by inserting $1,2,\dots,n$ one at a time: element $j$ starts a new
cycle or follows one of the $j-1$ previous elements), which splits. By (a), $c(n,k)^2\ge c(n,k-1)c(n,k+1)$. A positive log-concave
sequence is unimodal: the ratios $c(n,k+1)/c(n,k)$ are non-increasing, so they are $\ge1$ up to some $k$ and $\le1$ after.

\tools Tome~0, Chapter~3 (polynomials, Rolle's theorem), Chapter~8 (permutations and cycles).
\leads Newton stated (a) in 1707 without proof; the proof above is in Hardy, Littlewood and Pólya's \emph{Inequalities}. Log-concavity of
combinatorial sequences is a large subject: Heilmann and Lieb proved that matching numbers are real-rooted, and Huh and
collaborators proved Rota's and Mason's log-concavity conjectures with Hodge theory (2012--2020).
\end{solution}

\begin{solution}{pb:bonferroni}
\idea Look at one outcome at a time. If it lies in $r$ of the events, the truncated alternating sum counts it $1-(-1)^t\binom{r-1}t$ times,
which is above or below $1$ depending on the parity of $t$.

\writeup (a) For an outcome $\omega$ in exactly $r\ge1$ events, the indicator of the union is $1$, and $\omega$ contributes $\binom rj$ to $S_j$ (it lies in
$\binom rj$ of the intersections of $j$ events). By induction on $t$ from Pascal's rule,
$\sum_{j=1}^t(-1)^{j-1}\binom rj=1-(-1)^t\binom{r-1}t$. This is $\ge1$ for $t$ odd and $\le1$ for $t$ even, and for $r=0$ every term is $0$.
Taking expectations of these pointwise inequalities gives (a).

(b) Let $A_i$ be the event that $i$ is a fixed point. An intersection of $j$ of them has probability $\frac{(n-j)!}{n!}$, so $S_j=\binom nj\frac{(n-j)!}{n!}=\frac1{j!}$.
Then $\frac{D_n}{n!}=1-P(\bigcup A_i)$ lies between $\sum_{j=0}^t\frac{(-1)^j}{j!}$ and $\sum_{j=0}^{t+1}\frac{(-1)^j}{j!}$, which differ by $\frac1{(t+1)!}$ (for $t<n$;
for $t=n$ it is exact inclusion--exclusion).

\tools Tome~0, Chapter~8 (inclusion--exclusion, derangements).
\leads Brun's sieve (1915) is a clever truncation of inclusion--exclusion over primes, using Bonferroni at a well-chosen $t$; it gave
the first upper bound for twin primes (Tome~5, Chapter~8).
\end{solution}

\begin{solution}{pb:primitive}
\idea Attach to each $a$ the integers whose ``beginning'', in the factorisation sorted by primes, is exactly $a$. Primitivity makes these
sets disjoint, so their densities add up to at most $1$, and by Mertens the density attached to $a$ is about $\frac1{a\log a}$ or more.

\writeup (a) For $N=am\in M_a$, all primes of $m$ are $\ge P(a)$ while all primes of $a$ are $\le P(a)$. So $a$ is obtained from the prime
factorisation of $N$, listed in increasing order with multiplicity, by keeping an initial segment (all primes $<P(a)$ with their full
multiplicity, and some copies of $P(a)$). Two initial segments of the same list are comparable for divisibility. So if $N\in M_a\cap M_{a'}$
then $a\mid a'$ or $a'\mid a$; for a primitive set, $a=a'$.

(b) $N\in M_a$ if and only if $a\mid N$ and $N/a$ has no prime factor $<P(a)$. The second condition is periodic modulo $Q=\prod_{p<P(a)}p$ and holds
for a proportion $\prod_{p<P(a)}(1-\frac1p)$ of residues. So $M_a$ has density $\frac1a\prod_{p<P(a)}(1-\frac1p)$.

(c) By (a), for any finite subset $A'$ of $A$ the union of the $M_a$, $a\in A'$, has density $\sum_{a\in A'}d(M_a)\le1$ (densities of finitely many
disjoint periodic sets add). By Mertens's theorem (Tome~5, Theorem~1.6), there is $c>0$ with $\prod_{p<y}(1-\frac1p)\ge\frac c{\log y}$ for all $y\ge2$
(the asymptotic formula handles large $y$, and finitely many small $y$ are positive). As $P(a)\le a$,
\[
\sum_{a\in A'}\frac c{a\log a}\le\sum_{a\in A'}\frac1a\prod_{p<P(a)}\Bigl(1-\frac1p\Bigr)\le1,
\]
so $\sum_{a\in A}\frac1{a\log a}\le\frac1c$.

\tools Tome~5, Theorem~1.6 (Mertens); Tome~2, Chapter~2 (periodic sets and Chinese remainders); Tome~1, Chapter~2 (unique
factorisation).
\leads For the primes, $\sum_p\frac1{p\log p}=1.6366\ldots$, and Erdős conjectured in 1988 that no primitive set does better. Jared Duker
Lichtman proved it in 2022. Without the $\log a$, a primitive subset of $[1,x]$ can have $\sum\frac1a$ as large as $c\frac{\log x}{\sqrt{\log\log x}}$
(Besicovitch, Behrend).
\end{solution}

\begin{solution}{pb:ruzsa}
\idea An injection proves an inequality between sizes. For the covering, a maximal packing of translates is automatically a
covering by the doubled sets.

\writeup (a) For each $x\in B-C$ choose $b(x)\in B$, $c(x)\in C$ with $x=b(x)-c(x)$, and define $\phi:A\times(B-C)\to(A-B)\times(A-C)$ by
$\phi(a,x)=(a-b(x),a-c(x))$. From $\phi(a,x)=(u,v)$ we recover $x=v-u$ and then $a=u+b(x)$, so $\phi$ is injective. For the consequence, take
$B=C=-A$: then $|A|\,|A-A|\le|A+A|^2\le K^2|A|^2$.

(b) Let $X\subseteq B$ be maximal such that the translates $x+A$, $x\in X$, are pairwise disjoint. They lie in $A+B$, so $|X||A|\le|A+B|\le K|A|$ and
$|X|\le K$. For $b\in B$, maximality gives $x\in X$ with $(b+A)\cap(x+A)\ne\varnothing$ (if $b\notin X$; if $b\in X$ take $x=b$), so $b\in x+A-A$.

\tools Tome~2, Chapter~1 (abelian groups); Tome~0, Chapter~8 (injections).
\leads These are the first two steps of additive combinatorics. With Plünnecke's inequality ($|nA-mA|\le K^{n+m}|A|$) they lead to
Freiman's theorem: a set of integers with $|A+A|\le K|A|$ is contained in a generalised arithmetic progression of bounded dimension and
size $O_K(|A|)$. Gowers used them in his proof of Szemerédi's theorem.
\end{solution}

\begin{solution}{pb:divisorbound}
\idea $d(n)/n^\alpha$ is multiplicative, so its maximum is a product of local maxima, one for each prime; only the primes below
$2^{1/\alpha}$ can do better than $e=0$, and each local factor is unimodal in $e$.

\writeup (a) For $n=\prod p^{e_p}$, $\frac{d(n)}{n^\alpha}=\prod_pf_p(e_p)$ with $f_p(e)=\frac{e+1}{p^{e\alpha}}$, $f_p(0)=1$. For $p^\alpha\ge2$, $f_p(e)\le\frac{e+1}{2^e}\le1$,
so the maximum over $e$ is $1$; for each of the finitely many $p<2^{1/\alpha}$ the maximum over $e$ exists since $f_p(e)\to0$. Choosing each $e_p$ at
a maximiser shows that the product of local maxima is attained, and no $n$ exceeds it.

(b) The smallest maximiser $n_\alpha$ takes, for each $p$, the smallest maximiser of $f_p$. Now
\[
\frac{f_p(e+1)}{f_p(e)}=\frac{e+2}{e+1}\,p^{-\alpha},
\]
and $\frac{e+2}{e+1}$ decreases with $e$. So $f_p$ increases strictly as long as $\frac{e+2}{e+1}>p^\alpha$ and then never increases again: its
smallest maximiser is the least $e$ with $\frac{e+2}{e+1}\le p^\alpha$. When $p$ grows, $p^\alpha$ grows, the condition becomes easier and
this least $e$ can only decrease: $e_2\ge e_3\ge e_5\ge\cdots$. Finally $e_p\ge1$ if and only if the condition fails at $e=0$, that is $2>p^\alpha$.

(c) $\alpha=\frac13$: the primes with $p^{1/3}<2$ are those below $8$. The thresholds $\frac{e+2}{e+1}$ are $2,\frac32,\frac43,\frac54,\dots$, and
$2^{1/3}=1.260$, $3^{1/3}=1.442$, $5^{1/3}=1.710$, $7^{1/3}=1.913$. So $e_2=3$ (first threshold below $1.260$ is $\frac54$), $e_3=2$, $e_5=e_7=1$, and
$n_{1/3}=2^3\cdot3^2\cdot5\cdot7=2520$, with $d=48$ and $\frac{48}{2520^{1/3}}=3.527\ldots<3.53$.
For $\alpha=\frac14$ the primes with $p<16$ count, with $2^{1/4}=1.189$, $3^{1/4}=1.316$, $5^{1/4}=1.495$, $7^{1/4}=1.627$ and
$11^{1/4},13^{1/4}\in\left]1.5,2\right[$: so $e_2=5$ (the threshold $\frac76=1.167$), $e_3=3$ (the threshold $\frac54$), $e_5=2$ (the
threshold $\frac43$), $e_7=e_{11}=e_{13}=1$ (the threshold $\frac32$), and $n_{1/4}=2^5\cdot3^3\cdot5^2\cdot7\cdot11\cdot13=21\,621\,600$, with
$d(n_{1/4})=6\cdot4\cdot3\cdot2\cdot2\cdot2=576$ and $\frac{576}{n_{1/4}^{1/4}}=8.446\ldots$

\tools Tome~1, Chapter~4 (multiplicative functions).
\leads Wigert (1907) found the true order of the maximum: $\log d(n)\le(\log2+o(1))\frac{\log n}{\log\log n}$, with equality infinitely often.
Ramanujan studied the record holders, the \emph{highly composite numbers}, such as $12$, $2520$ and $5040$, in a long paper of 1915.
\end{solution}

\section*{Map of the chapter}

\begin{center}
\small
\adjustbox{max width=\linewidth}{\begin{tabular}{lll}
Problem & Topic & Tomes used\\\hline
\ref{pb:distances} & distinct and unit distances & 0, 3\\
\ref{pb:incidences} & point-line incidences by Cauchy--Schwarz & 0\\
\ref{pb:szemereditrotter} & Szemerédi--Trotter, unit distances & 3\\
\ref{pb:sumproduct} & Elekes's sum-product bound & 3\\
\ref{pb:energy} & energy, the multiplication table & 0, 1, 5\\
\ref{pb:triangleseigen} & triangles and eigenvalues & 0, 11\\
\ref{pb:goodman} & monochromatic triangles & 0, 3\\
\ref{pb:walks} & walks, $4$-cycles, Paley graphs & 0, 11\\
\ref{pb:treebound} & spanning trees and AM--GM & 0, 11\\
\ref{pb:bollobas} & set-pairs inequality, $\tau$-critical graphs & 3, 3\\
\ref{pb:littlewoodofford} & anti-concentration of random sums & 0, 3\\
\ref{pb:kleitman} & Harris--Kleitman correlation & 0\\
\ref{pb:sauer} & VC-dimension & 0\\
\ref{pb:cubeiso} & edge isoperimetry in the cube & 0\\
\ref{pb:erdosselfridge} & Maker--Breaker games & 0, 3\\
\ref{pb:shearer} & entropy, Loomis--Whitney & 3\\
\ref{pb:chebyshevsum} & Chebyshev's sum inequality & 3\\
\ref{pb:newton} & Newton's inequalities, log-concavity & 0\\
\ref{pb:bonferroni} & Bonferroni's inequalities & 0\\
\ref{pb:primitive} & primitive sets & 1, 2, 5\\
\ref{pb:ruzsa} & Ruzsa's inequalities & 0, 2\\
\ref{pb:divisorbound} & the size of $d(n)$ & 1\\
\end{tabular}}
\end{center}

\section*{Sources and further reading}

\begin{itemize}
\item M.~Aigner and G.~M.~Ziegler, \emph{Proofs from THE BOOK}, 6th ed., Springer, 2018 (the epigraph is quoted in the preface; chapters
on the Littlewood--Offord problem, the crossing lemma and Sperner's theorem).
\item P.~Erdős, On sets of distances of $n$ points, \emph{American Mathematical Monthly} 53 (1946), 248--250.
\item L.~A.~Székely, Crossing numbers and hard Erdős problems in discrete geometry, \emph{Combinatorics, Probability and Computing} 6
(1997), 353--358; E.~Szemerédi and W.~T.~Trotter, Extremal problems in discrete geometry, \emph{Combinatorica} 3 (1983), 381--392.
\item G.~Elekes and M.~Sharir, Incidences in three dimensions and distinct distances in the plane, \emph{Combinatorics, Probability and
Computing} 20 (2011), 571--608; L.~Guth and N.~H.~Katz, On the Erdős distinct distances problem in the plane, \emph{Annals of
Mathematics} 181 (2015), 155--190.
\item G.~Elekes, On the number of sums and products, \emph{Acta Arithmetica} 81 (1997), 365--367; J.~Solymosi, Bounding multiplicative
energy by the sumset, \emph{Advances in Mathematics} 222 (2009), 402--408.
\item K.~Ford, The distribution of integers with a divisor in a given interval, \emph{Annals of Mathematics} 168 (2008), 367--433.
\item A.~W.~Goodman, On sets of acquaintances and strangers at any party, \emph{American Mathematical Monthly} 66 (1959), 778--783.
\item F.~R.~K.~Chung, R.~L.~Graham and R.~M.~Wilson, Quasi-random graphs, \emph{Combinatorica} 9 (1989), 345--362.
\item B.~Bollobás, On generalized graphs, \emph{Acta Mathematica Academiae Scientiarum Hungaricae} 16 (1965), 447--452.
\item P.~Erdős, On a lemma of Littlewood and Offord, \emph{Bulletin of the AMS} 51 (1945), 898--902.
\item T.~E.~Harris, A lower bound for the critical probability in a certain percolation process, \emph{Proceedings of the Cambridge
Philosophical Society} 56 (1960), 13--20; D.~J.~Kleitman, Families of non-disjoint subsets, \emph{Journal of Combinatorial Theory} 1
(1966), 153--155.
\item N.~Sauer, On the density of families of sets, \emph{Journal of Combinatorial Theory A} 13 (1972), 145--147.
\item P.~Erdős and J.~L.~Selfridge, On a combinatorial game, \emph{Journal of Combinatorial Theory A} 14 (1973), 298--301.
\item F.~R.~K.~Chung, R.~L.~Graham, P.~Frankl and J.~B.~Shearer, Some intersection theorems for ordered sets and graphs, \emph{Journal of
Combinatorial Theory A} 43 (1986), 23--37; L.~H.~Loomis and H.~Whitney, An inequality related to the isoperimetric inequality,
\emph{Bulletin of the AMS} 55 (1949), 961--962.
\item G.~H.~Hardy, J.~E.~Littlewood and G.~Pólya, \emph{Inequalities}, 2nd ed., Cambridge, 1952 (Chebyshev's and Newton's inequalities).
\item P.~Erdős, Note on sequences of integers no one of which is divisible by any other, \emph{Journal of the London Mathematical
Society} 10 (1935), 126--128; J.~D.~Lichtman, A proof of the Erdős primitive set conjecture, \emph{Forum of Mathematics, Pi} 11 (2023), e18.
\item I.~Z.~Ruzsa, Sums of finite sets, in \emph{Number Theory: New York Seminar 1991--1995}, Springer, 1996; T.~Tao and V.~Vu,
\emph{Additive Combinatorics}, Cambridge, 2006.
\item S.~Ramanujan, Highly composite numbers, \emph{Proceedings of the London Mathematical Society} (2) 14 (1915), 347--409.
\end{itemize}
